\chapter{Gestão e Compartilhamento de Dados de Pesquisa}\label{chp:FUNDAMENTACAO}

Este capítulo discute a problemática que motiva a proposta do DataRural. A questão central é a gestão dos dados produzidos em atividades de pesquisa: como organizá-los, documentá-los, preservá-los e compartilhá-los de forma que permaneçam úteis para além do projeto ou da pessoa que os gerou. Para situar essa discussão, são apresentados os conceitos de dataset e metadados, os desafios enfrentados por instituições que produzem dados sem infraestrutura dedicada, os princípios que orientam boas práticas de disponibilização, as implicações do licenciamento e da privacidade, e o papel do versionamento na preservação do histórico. O capítulo também aborda os tipos de plataformas que procuram responder a esses desafios e as particularidades do formato CSV como veículo de dados tabulares.

\section{Dados de pesquisa, datasets e metadados}

Universidades produzem dados em diferentes contextos: experimentos de laboratório, levantamentos de campo, registros administrativos, indicadores acadêmicos e projetos de extensão. Esses registros podem apoiar novos estudos, subsidiar decisões institucionais ou servir como evidência para resultados publicados. Entretanto, a simples existência de um arquivo não garante que seu conteúdo será compreendido, localizado ou reutilizado por outras pessoas.

Um \textit{dataset}, ou conjunto de dados, é uma coleção organizada de registros relacionados. Em uma representação tabular, cada linha geralmente corresponde a uma observação e as colunas representam atributos. Quando um dataset é publicado sem informações sobre sua origem, finalidade, unidade responsável, período coberto ou procedimentos de coleta, o custo de interpretá-lo pode superar o custo de produzir novos dados. Esse problema é agravado quando os dados passam de uma pessoa para outra ou quando o responsável original não está mais disponível para esclarecer dúvidas.

Essas informações complementares são denominadas metadados. São dados que descrevem outros dados e oferecem elementos para sua localização, interpretação, administração e reutilização \cite{riley2017metadata}. Metadados descritivos incluem título, autoria, resumo e palavras-chave; metadados administrativos registram licença, responsável e condições de acesso; e metadados estruturais informam como os arquivos, tabelas e variáveis se relacionam \cite{riley2017metadata}. A ausência de qualquer dessas camadas pode tornar um conjunto tecnicamente acessível, mas praticamente inutilizável.

\section{Gestão de dados de pesquisa}\label{sec:rdm}

A gestão de dados de pesquisa, conhecida como \textit{research data management} (RDM), reúne as práticas aplicadas ao longo do ciclo de vida dos dados produzidos em atividades de investigação \cite{higgins2008dcc}. Esse ciclo inclui planejamento, coleta, organização, documentação, armazenamento, controle de acesso, preservação, compartilhamento e, eventualmente, descarte. Uma gestão adequada reduz a dependência do conhecimento individual de quem produziu os arquivos e aumenta a possibilidade de utilização futura por outras pessoas.

A discussão sobre gestão e compartilhamento envolve tanto as possibilidades de reutilização quanto os incentivos e as responsabilidades dos produtores dos dados. \citeonline{borgman2012conundrum} distingue motivações como a verificação de pesquisas, o acesso aos resultados do financiamento público, a formulação de novas perguntas e o avanço científico. Essas finalidades não pressupõem práticas idênticas entre as disciplinas, mas evidenciam a necessidade de preservar o contexto dos registros e de explicitar as condições de compartilhamento.

\subsection{Desafios da gestão de dados}

Apesar do reconhecimento crescente, a gestão de dados enfrenta desafios práticos que dificultam sua adoção. Em um estudo de dez anos no \textit{Center for Embedded Networked Sensing}, \citeonline{wallis2013sharing} observaram que o compartilhamento ocorria principalmente por trocas interpessoais e dependia de condições como o reconhecimento dos produtores e a prioridade de publicação. Embora os resultados se refiram ao centro estudado, eles ajudam a compreender por que disponibilizar infraestrutura não basta para estabelecer práticas de compartilhamento.

O primeiro desafio é a fragmentação. Os dados permanecem distribuídos entre computadores pessoais, pastas compartilhadas, serviços de armazenamento em nuvem e sistemas internos sem integração. Quando um pesquisador se aposenta, muda de instituição ou simplesmente conclui um projeto, os arquivos produzidos podem tornar-se inacessíveis ou incompreensíveis para outras pessoas. A continuidade do valor dos dados depende de registros que não se percam junto com a memória individual de seus produtores.

O segundo desafio é a perda de contexto. Mesmo quando um arquivo pode ser localizado, nem sempre estão disponíveis informações suficientes para esclarecer sua origem, o período representado, a unidade responsável ou as condições de utilização. Sem metadados adequados, dois pesquisadores podem utilizar o mesmo arquivo e produzir interpretações divergentes sem perceber que estão tratando de versões diferentes ou de recortes distintos dos mesmos dados.

O terceiro desafio é a reprodutibilidade. Alterações realizadas diretamente sobre um mesmo documento podem apagar estados anteriores e impedir a verificação de análises passadas. Se não existe um mecanismo que preserve cada estado do conjunto, torna-se impossível determinar qual material sustentou um resultado previamente publicado \cite{klump2021versioning}.

O quarto desafio envolve as barreiras à reutilização. Um dataset pode estar acessível publicamente e, ainda assim, não conter informações suficientes sobre o método de coleta, as unidades de medida, os valores ausentes ou as limitações conhecidas. Nesses casos, o esforço para compreender e validar os dados pode anular a vantagem do compartilhamento.

Os problemas discutidos na literatura orientam a análise das necessidades do DataRural, mas não constituem, por si só, um diagnóstico empírico da UFRRJ. Neste trabalho, a fragmentação, a perda de contexto e as barreiras à reutilização são tomadas como riscos a enfrentar na concepção de uma plataforma institucional. A proposta procura reduzir a dependência de arquivos dispersos e de esclarecimentos individuais, reunindo publicação, documentação e controle de acesso.

\subsection{Privacidade, ética e dados sensíveis}

Nem todos os dados produzidos em atividades de pesquisa podem ou devem ser disponibilizados livremente. Pesquisas que envolvem dados pessoais, informações de saúde, registros administrativos ou materiais protegidos por sigilo institucional exigem cuidados específicos que antecedem qualquer decisão de publicação.

No Brasil, a Lei Geral de Proteção de Dados Pessoais (LGPD, Lei nº~13.709/2018) estabelece regras para o tratamento de dados pessoais, incluindo a coleta, o armazenamento e o compartilhamento dessas informações \cite{brasil2018lgpd}. A lei exige que o tratamento observe princípios como finalidade, adequação, necessidade e transparência. Pesquisas acadêmicas possuem previsões específicas, mas não estão isentas das obrigações de proteção.

Além das exigências legais, existem questões éticas. Dados de populações vulneráveis, comunidades tradicionais ou grupos sub-representados podem conter vieses que, se não forem documentados, podem levar à perpetuação de desigualdades em análises derivadas. A documentação das características e limitações dos dados contribui para que seus reutilizadores avaliem a adequação do material ao contexto pretendido.

Essas restrições tornam evidente que uma plataforma de gestão de dados não pode se limitar à publicação aberta. É necessário que o responsável possa definir o nível de acesso de cada conjunto e que o sistema ofereça mecanismos para que essas condições sejam respeitadas. A abertura total não é desejável em todos os casos, e a restrição de acesso não deve ser confundida com falta de transparência.

\section{Princípios FAIR}\label{sec:fair}

Diante dos desafios apresentados na seção anterior, foram propostos princípios que orientam a descrição e a disponibilização de dados de forma a favorecer sua descoberta, acesso e reutilização. Os princípios FAIR, publicados por \citeonline{wilkinson2016fair}, estabelecem que dados e outros objetos digitais devem ser:

\begin{compactitem}
    \item Localizáveis (\textit{findable}): os dados devem ser descritos com metadados ricos e receber identificadores persistentes que facilitem sua descoberta por pessoas e por sistemas computacionais;
    \item Acessíveis (\textit{accessible}): uma vez localizados, deve existir um procedimento claro para recuperar os dados, ainda que esse procedimento envolva autenticação ou autorização;
    \item Interoperáveis (\textit{interoperable}): os dados e seus metadados devem utilizar vocabulários e formatos que permitam a interpretação e o cruzamento com outros conjuntos;
    \item Reutilizáveis (\textit{reusable}): os dados devem ser acompanhados de informações de proveniência, licença e contexto suficientes para que terceiros possam avaliá-los e utilizá-los de forma apropriada.
\end{compactitem}

É importante observar que os princípios FAIR não exigem que todos os dados sejam abertos. A acessibilidade prevista pelo segundo princípio reconhece que o acesso pode depender de procedimentos de autenticação e autorização \cite{wilkinson2016fair}. Um dataset pode seguir os princípios FAIR e, ainda assim, ter acesso restrito a um grupo de pesquisa ou a usuários autorizados. O que se espera é que as condições de acesso estejam documentadas e que exista um caminho definido para solicitá-lo.

A rastreabilidade complementa os princípios FAIR ao permitir identificar a origem e o histórico de um dado, incluindo as transformações ou versões pelas quais passou. Sem essa informação, duas análises que utilizam arquivos com o mesmo nome podem produzir resultados diferentes sem que a causa seja facilmente identificada.

\section{Dados abertos e licenciamento}

A reutilização de dados depende não apenas de sua disponibilidade técnica, mas também da clareza das condições de uso. Ao analisar 56 recursos de dados científicos no contexto biomédico, \citeonline{carbon2019licensing} identificaram barreiras decorrentes de licenças ausentes, restritivas ou incompatíveis. O estudo distingue o acesso aos dados da possibilidade de reutilizá-los e redistribuí-los. Para o DataRural, essa distinção fundamenta a necessidade de apresentar as condições de uso junto aos metadados, sem tratar o acesso público como autorização irrestrita.

O uso de instrumentos padronizados pode reduzir ambiguidades, embora não elimine todas as restrições à reutilização \cite{carbon2019licensing}. As licenças CC BY e CC BY-SA estabelecem condições de atribuição e, no segundo caso, de compartilhamento das adaptações sob a mesma licença \cite{creativecommonslicenses}. O CC0 é um instrumento de dedicação ao domínio público, e não uma licença de atribuição \cite{creativecommonscczero}. A Open Database License (ODbL), por sua vez, estabelece condições para o uso e o compartilhamento de bases de dados \cite{opendatacommonsodbl}. A escolha de um instrumento deve considerar os direitos sobre o material e não afasta cuidados com privacidade ou sigilo.

Assim, abertura e controle de acesso não são conceitos incompatíveis. Uma mesma plataforma pode reunir conjuntos abertos, conjuntos públicos com condições específicas e conjuntos privados destinados a uma equipe. O aspecto central é que o nível de acesso e as condições de reutilização sejam apresentados de maneira explícita, evitando a ambiguidade que dificulta o reaproveitamento.

\section{Versionamento de dados}\label{sec:versionamento}

Datasets de pesquisa são frequentemente objetos em evolução. Novos registros podem ser adicionados à medida que a coleta avança, correções podem alterar valores previamente publicados e recodificações podem modificar a estrutura das variáveis. Quando essas mudanças são aplicadas diretamente sobre o mesmo arquivo, sem preservar os estados anteriores, torna-se impossível determinar qual versão dos dados sustentou um resultado previamente divulgado.

O versionamento responde a esse problema ao preservar estados sucessivos de um recurso ao longo do tempo. Em vez de substituir definitivamente um arquivo a cada alteração, o sistema cria uma nova versão e mantém a relação com as anteriores. Essa abordagem permite recuperar estados passados e informar com precisão qual versão foi utilizada em determinada análise.

A identificação precisa das versões permite relacionar os dados aos resultados que deles derivam. \citeonline{klump2021versioning}, a partir da análise de práticas de 33 organizações, propõem princípios que distinguem revisão, publicação de versões, granularidade, manifestação, proveniência e citação. Essa abordagem evidencia que registrar apenas um número de versão é insuficiente: também é necessário documentar a natureza das alterações e manter relações que permitam identificar o material utilizado em uma análise.

O versionamento de dados também exige decisões de armazenamento diferentes das adotadas para código-fonte. \citeonline{hoyt2024o3} observam que sistemas de controle de versão são otimizados para arquivos textuais e podem ser inadequados para arquivos muito grandes ou binários. Os autores apontam repositórios de dados como alternativas para armazenar e compartilhar esses materiais. No DataRural, a preservação dos arquivos de cada versão é a estratégia escolhida para manter estados anteriores dos conjuntos, sem depender da reconstrução por diferenças entre linhas.

\section{Plataformas de compartilhamento de dados}\label{sec:plataformas}

Os desafios descritos nas seções anteriores motivaram o desenvolvimento de plataformas que procuram apoiar diferentes etapas da gestão de dados. Essas plataformas assumem formatos distintos de acordo com o público atendido e a finalidade dos conjuntos publicados.

\citeonline{pampel2013making} discutem a diversidade dos repositórios de dados de pesquisa e apresentam o registro re3data.org como instrumento para sua descoberta. A descrição dessas infraestruturas envolve aspectos como escopo, acesso, licenças e identificadores persistentes. Para o presente trabalho, esses aspectos oferecem critérios de comparação mais precisos do que a simples presença de armazenamento e download.

A literatura também distingue repositórios de propósito geral de sistemas especializados em determinados domínios. Em uma análise voltada a dados de pesquisa matemática, \citeonline{conrad2024repositories} incluem Zenodo e Figshare entre as soluções de propósito geral e discutem recursos como identificação persistente e preservação. Essa distinção ajuda a situar o DataRural: seu recorte é institucional, enquanto as plataformas relacionadas atendem comunidades mais amplas. Independentemente do recorte, autoria, licença, versão, descrição e condições de acesso devem acompanhar o arquivo publicado.

A existência dessas plataformas não resolve, por si só, todos os problemas da gestão de dados. A qualidade dos metadados, a adesão dos pesquisadores e a manutenção institucional continuam sendo fatores determinantes. Ainda assim, uma solução tecnológica que centralize a publicação, a documentação e o controle de acesso pode reduzir a fragmentação e oferecer um ponto de referência para os dados produzidos em uma instituição. A análise detalhada de plataformas existentes é apresentada na seção de trabalhos relacionados do Capítulo~\ref{chp:PROPOSTA}.

\section{Formato CSV e dados tabulares}

O formato CSV (\textit{comma-separated values}) é amplamente utilizado como veículo de dados tabulares por sua simplicidade e compatibilidade \cite{shafranovich2005csv}. Entretanto, essa mesma simplicidade introduz problemas relevantes para a gestão de dados.

O CSV não define obrigatoriamente os tipos das colunas, não mantém relacionamentos entre tabelas e não possui um esquema embutido capaz de descrever cada variável \cite{w3c2015csvw}. Datas, números decimais, valores ausentes, delimitadores e codificações de caracteres podem ser interpretados de formas diferentes por aplicações distintas. Dois arquivos identificados como CSV podem exigir configurações completamente diferentes durante a importação.

Essas limitações tornam os metadados especialmente importantes quando o CSV é o formato adotado. O nome de uma coluna nem sempre explica sua unidade de medida, os valores permitidos ou o procedimento utilizado na coleta. Um arquivo pode ser amplamente compatível do ponto de vista técnico, mas sua reutilização efetiva depende de documentação externa e de convenções consistentes. Sem essa documentação, a simplicidade do formato torna-se uma armadilha: o dado é acessível, mas não compreensível.

Uma prévia tabular apresenta as primeiras linhas e colunas de um arquivo antes do download completo. Essa visualização auxilia a inspeção rápida e permite reconhecer a estrutura do conjunto. Esse recurso, entretanto, não substitui uma análise de qualidade, uma vez que problemas podem existir em linhas não apresentadas ou depender do significado científico dos valores.
